% ============================================================
%  CAPÍTULO 13 – Plan de Ejecución del Proyecto
% ============================================================
\chapter{Plan de Ejecución del Proyecto}

\section{Permisos y autorizaciones}

Para el desarrollo de FitLabs, los permisos necesarios son los siguientes:

\begin{itemize}
  \item \textbf{Permisos en Android}: La aplicación requiere acceso a Internet (\texttt{INTERNET}) declarado en el \texttt{AndroidManifest.xml}, necesario para la comunicación con Supabase y la API de ejercicios.
  \item \textbf{Acceso a Supabase}: Las credenciales de acceso al proyecto Supabase (URL y \texttt{anon key}) se gestionan de forma segura y no se incluyen en el repositorio público de GitHub.
  \item \textbf{Claves de API externas}: La API key de acceso al catálogo de ejercicios se almacena como constante en el código fuente, protegida mediante el plan de acceso del repositorio privado.
\end{itemize}

\section{Procedimientos de ejecución}

\subsection{Flujo de trabajo con Git}

El equipo siguió el siguiente flujo de trabajo:

\begin{enumerate}
  \item Cada desarrollador trabajó en su rama personal (\texttt{Luis} -- José Luis Sánchez González -- y \texttt{Carlos} -- Carlos Fraidias del Valle --).
  \item Al completar una funcionalidad, se realizaba un \textit{commit} descriptivo con el prefijo del módulo (e.g., \texttt{feat(clientes): añadir buscador por nombre}).
  \item Se revisaba el código antes de integrar los cambios a \texttt{main} mediante \textit{merge} manual.
\end{enumerate}

\subsection{Proceso de despliegue en Android}

\begin{enumerate}
  \item Compilar el APK en modo \textit{release}: \texttt{flutter build apk}.
  \item Firmar el APK con un \textit{keystore} para distribución.
  \item Instalar en el dispositivo de prueba mediante \texttt{adb install} o transferencia directa.
\end{enumerate}

\section{Identificación de riesgos y medidas preventivas}

\begin{longtable}{@{}p{5cm}p{3cm}p{6cm}@{}}
\toprule
\textbf{Riesgo} & \textbf{Probabilidad} & \textbf{Medida preventiva} \\
\midrule
\endhead
Cambio en la API de ejercicios (fin de servicio o cambio de términos) & Media & Abstraer las llamadas en un repositorio (\texttt{ExerciseRepository}) para facilitar el cambio de proveedor \\
Pérdida de código (fallo de disco, accidente) & Baja & Repositorio en GitHub con \textit{push} frecuente; copia de seguridad del proyecto en Google Drive \\
Superación del límite del plan Free de Supabase & Baja & Monitorización de uso en el dashboard; posibilidad de migrar a plan Pro \\
Incompatibilidades entre versiones de Flutter & Baja & Fijar la versión del SDK en \texttt{pubspec.yaml} \\
Falta de tiempo para completar todos los módulos & Media & Priorización de módulos por valor funcional; módulos secundarios postergados a Vías Futuras \\
\bottomrule
\caption{Tabla de riesgos identificados y medidas preventivas.}
\end{longtable}

\section{Estrategia de integración y control de versiones}

Para minimizar incidencias en la rama principal, se siguió una estrategia de integración progresiva:

\begin{itemize}
  \item \textbf{Ramas por funcionalidad}: cada bloque funcional (clientes, rutinas, calendario, mensajes) se desarrolló en una rama independiente.
  \item \textbf{Commits atómicos}: se evitó mezclar cambios de UI con cambios de base de datos en un mismo commit.
  \item \textbf{Merge controlado}: antes de fusionar a \texttt{main}, se validaba compilación en local y ejecución de casos críticos.
  \item \textbf{Etiquetado de hitos}: se marcaron versiones internas para entregas parciales y revisión con tutor.
\end{itemize}

Esta disciplina permitió reducir conflictos de integración y facilitó la trazabilidad de incidencias por módulo.

\section{Checklist de entrega técnica}

Antes de cada entrega formal se aplicó una lista de verificación:

\begin{longtable}{@{}p{8.5cm}p{2.5cm}p{3cm}@{}}
\toprule
\textbf{Elemento de control} & \textbf{Estado} & \textbf{Responsable} \\
\midrule
\endhead
Compilación Flutter sin errores críticos & Completado & Equipo \\
Conexión a Supabase validada en entorno de pruebas & Completado & Equipo \\
Migraciones SQL verificadas en proyecto remoto & Completado & Equipo \\
Revisión de credenciales fuera del repositorio & Completado & Equipo \\
Revisión de capturas y figuras de memoria & Completado & Equipo \\
Compilación final de la memoria en PDF & Completado & Equipo \\
\bottomrule
\caption{Checklist de control previo a entrega.}
\end{longtable}

\section{Plan de continuidad y mantenimiento inicial}

Aunque el objetivo principal es académico, se definió un plan básico de continuidad para facilitar la evolución del producto tras la entrega:

\begin{enumerate}
  \item \textbf{Mantenimiento correctivo}: resolución prioritaria de bugs de autenticación, calendario y mensajería.
  \item \textbf{Mantenimiento adaptativo}: actualización de dependencias de Flutter/Supabase cuando exista incompatibilidad relevante.
  \item \textbf{Mantenimiento evolutivo}: incorporación de funcionalidades propuestas en Vías Futuras según valor para el usuario final.
\end{enumerate}

Este plan garantiza que el proyecto no quede como un entregable cerrado, sino como una base sostenible para una versión posterior.
